\section{Higher heights as modular counting problems}\label{fg:sec:counting}
The next height test is a residue of a count associated with the incidence
presentation. It depends on more than the number of satisfying assignments.
Nevertheless, every fixed threshold at least two is modular-counting complete.
We prove the coefficient identities first, then establish membership for
arbitrary arithmetic circuits and give the reductions.

Throughout this section $p$ is fixed. For a homogeneous integer form $H$
of degree $u$ in $u$ variables, write
\[
 B_t(H)=[(x_1\cdots x_u)^{t-1}]H^{t-1}\qquad(t\ge1).
\]
When $H$ is a chosen lift of a defining equation, this is the logarithm
numerator $b_t$ of \eqref{fg:eq:stienstra}. The notation also applies to
forms whose zero schemes are singular or reducible; no formal group is
attached to such an intermediate form.

\subsection{Replica coefficients and the second invariant}
Let $\varphi$ be a width-three CNF with $S$ satisfying assignments, and use
the integer incidence polynomials $P,Z$ of \eqref{eq:incidence}. Every
declared variable occurs in an incidence, after the prescribed
tautological relations and, if necessary, a forced variable are added.
Both $P$ and $Z$ have degree $b$ in $b$ variables, and $[Z]P=S$.
A \emph{local configuration} chooses one allowed term in every factor
$P_f$ and one of the two terms in every factor $Q_i$. Distinct choices
remain distinct even when they produce the same monomial.

Regard \eqref{eq:normalization} as an integer polynomial $G_\varphi$ in
$u=b+p-1$ variables. We likewise use the displayed integer lift of
\eqref{eq:stabilizer} when its input is an integer form.

\begin{lemma}[Replica numerators]\label{fg:lem:replica}
Let $G\in\Z[x_1,\ldots,x_u]$ be homogeneous of degree $u\ge1$, and let
$Q$ be a positive multiple of $p$ with $Q\ge u+2$.
For every $1\le t\le Q$,
\begin{equation}\label{fg:eq:stabilized-numerator}
 b_t(F_G)=B_t(G).
\end{equation}
For $i\ge1$, set
\[
 k_i=\frac{p^i-1}{p-1},\qquad
 A_i(\varphi)=[Z^{k_i}]P^{k_i},\qquad
 C_i=(-1)^{k_i}\frac{(p^i-1)!}{(k_i!)^{p-1}}.
\]
If $Q\ge p^i$, then
\begin{equation}\label{fg:eq:replica-numerator}
 b_{p^i}(F_{G_\varphi})=C_iA_i(\varphi)
\end{equation}
over the integers. Moreover
\begin{equation}\label{fg:eq:replica-residue}
 v_p(C_i)=0,\qquad C_i\equiv1\pmod p,\qquad
 A_i(\varphi)\equiv S^i\pmod p.
\end{equation}
For each fixed $p,i$, the function $A_i$ belongs to $\#\mathrm P$.
\end{lemma}

\begin{proof}
Put $K=t-1$. In a term contributing to
$(x_1\cdots x_myz)^K$ in $F_G^K$, the summand $-x_1z^Q$ cannot occur,
since $Q>K$. The only remaining source of $z$ is $yzA_G$, whose
$z$-exponent is one. Thus all $K$ factors select $yzA_G$. The padding
coordinate $x_{u+1}$ exists because $Q\ge u+2$. Its exponent is one in
$G\prod_{j>u}x_j$ and zero in $x_1^{Q-1}$, so its target exponent forces
the first summand in every factor. Removing the padding gives
\eqref{fg:eq:stabilized-numerator}. This also treats the endpoint $Q=t$.

Now take $t=p^i$. Each tag variable occurs in one summand of
$G_\varphi$, with exponent $p-1$. Its target exponent $p^i-1$ forces
exactly $k_i$ copies of that summand. The multinomial coefficient is
$(p^i-1)!/(k_i!)^{p-1}$. The first summand contributes
$(-1)^{k_i}P^{k_i}$ and the remaining summands contribute
$Z^{(p-2)k_i}$. Removing them leaves the coefficient
$[Z^{k_i}]P^{k_i}$, proving \eqref{fg:eq:replica-numerator}. At $p=2$
there is only the first tag summand, with the same conclusion.

For an integer $a\ge0$, let $s_p(a)$ be its base-$p$ digit sum. Counting
the multiples of each power of $p$ in $a!$ gives
$v_p(a!)=(a-s_p(a))/(p-1)$. The digits of $p^i-1$ are all $p-1$,
whereas those of $k_i$ are all one. Hence
\[
 v_p((p^i-1)!)=k_i-i,\qquad
 (p-1)v_p(k_i!)=k_i-i.
\]
This proves the unit assertion. The digitwise multinomial expansion
gives the factorial quotient modulo $p$ as
$((p-1)!)^i=(-1)^i$. For odd $p$, $k_i\equiv i\pmod2$; at $p=2$
the sign is immaterial. Therefore $C_i\equiv1$.

In characteristic $p$,
\[
 P^{k_i}=\prod_{j=0}^{i-1}P(w^{p^j}).
\]
For a contribution to $Z^{k_i}$, every exponent in the lowest block is
congruent to one modulo $p$, and is therefore at least one. The block
has total degree $b$ in $b$ variables, so all its exponents are exactly
one. Subtract this exponent vector and divide the remaining target by
$p$. Induction forces the selected monomial in every block to be $Z$.
Its coefficient is $S$ modulo $p$, proving $A_i\equiv S^i$.

Finally, guess $k_i$ ordered local configurations and accept exactly
when their monomial product is $Z^{k_i}$. Each clause factor has at most
eight choices and each variable factor two. The exponent test and the
witness length are polynomial because $p,i$, and hence $k_i$, are
fixed. This counts every coefficient contribution with its local-choice
multiplicity, proving the $\#\mathrm P$ assertion.
\end{proof}

The same degree argument shows that for every homogeneous degree-$u$
form $H$ in $u$ variables,
\[
 B_{p^i}(H)\equiv B_p(H)^i\pmod p.
\]
Indeed, split $H^{p^i-1}$ into Frobenius twists of $H^{p-1}$. In the
lowest block every exponent is congruent to $p-1$ and hence at least
$p-1$; its total degree $u(p-1)$ forces equality in every coordinate.
No individual-exponent bound is needed for this congruence.

\begin{theorem}[The second invariant as an orbit count]\label{fg:thm:orbit}
There is an explicit function $\mathcal C\in\#\mathrm P$ such that
\begin{equation}\label{fg:eq:orbit-count}
 A_2(\varphi)=S^2+p\mathcal C(\varphi).
\end{equation}
It counts the nontrivial orbits of the cyclic group of order $p$ on
balanced ordered collections of $p+1$ local configurations, rotating
the first $p$ configurations and fixing the last. Balance means that
their monomial product is $Z^{p+1}$.

For the original formula construction, or any admissible choice
$Q\ge p^2$, the second invariant on $S\equiv0\pmod p$ is
\[
 \frac{b_{p^2}}p\equiv\mathcal C(\varphi)\pmod p.
\]
Consequently
\begin{equation}\label{fg:eq:height-three-count}
 h(X_\varphi)\ge3
 \iff S\equiv\mathcal C(\varphi)\equiv0\pmod p
 \iff A_2(\varphi)\equiv0\pmod{p^2}.
\end{equation}
\end{theorem}

\begin{proof}
The set of balanced ordered collections has cardinality $A_2$, since
$k_2=p+1$. Its cyclic action preserves balance. Since $p$ is prime,
each orbit has size one or $p$, and a fixed record has its first $p$
configurations identical.

Every incidence coordinate appears in exactly one clause factor and
one variable factor. Thus its exponent in a local configuration is
zero, one, or two. If the common first monomial has exponent vector
$a$ and the last has exponent vector $c$, balance gives
\[
 pa_j+c_j=p+1,\qquad 0\le a_j,c_j\le2.
\]
The only possibility is $a_j=c_j=1$: $a_j=0$ would require
$c_j=p+1>2$, and $a_j=2$ would require $c_j=1-p<0$.
Both selected monomials are therefore $Z$. Lemma~\ref{lem:incidence}
gives $S$ local configurations for each, so there are $S^2$ fixed
records. Counting the remaining size-$p$ orbits proves
\eqref{fg:eq:orbit-count}.

To realize $\mathcal C$ itself in $\#\mathrm P$, encode each local
configuration injectively in a fixed-length string by ordering its
factor choices. Guess a balanced collection, reject fixed records,
and accept exactly when its encoding is strictly lexicographically
smaller than each of its nonidentity rotations. Every nontrivial orbit
has exactly one such record. The test is polynomial time for fixed $p$.
This is a canonical-representative count, not an assumed closure under
division by $p$.

The original choice of $Q$ is a power of $p$ with $Q-1\ge p$, hence
$Q\ge p^2$. Lemma~\ref{fg:lem:replica} gives $b_p=C_1S$ and
$b_{p^2}=C_2A_2$. If $p\mid S$, then
\[
 \frac{b_{p^2}}p
 =C_2\left(\frac{S^2}p+\mathcal C(\varphi)\right)
 \equiv\mathcal C(\varphi)\pmod p,
\]
using $p^2\mid S^2$ and $C_2\equiv1$. Apply
Corollary~\ref{fg:cor:second}. Finally, $p^2\mid A_2$ already implies
$S^2\equiv0\pmod p$ by \eqref{fg:eq:replica-residue}, so it forces
$p\mid S$. Equation~\eqref{fg:eq:orbit-count} then gives the remaining
equivalence in \eqref{fg:eq:height-three-count}.
\end{proof}

\begin{example}[Equivalent formulas with different second heights]
\label{fg:ex:same-function}
Work over $\F_2$ with the three declared Boolean variables
$x_0,x_1,x_2$. Consider
\begin{align*}
 \varphi_0={}&x_1\wedge x_2,\\
 \varphi_1={}&(\neg x_0\vee\neg x_1\vee x_2)
 \wedge(x_0\vee\neg x_1\vee x_2)\\
 &{}\wedge(\neg x_0\vee x_2)\wedge x_1.
\end{align*}
The unit clause forces $x_1=1$. The first two clauses of $\varphi_1$
then force $x_2=1$, and the third is redundant. Thus both formulas
compute the same Boolean function and have $S=2$.

Here $k_2=3$. We compute the coefficients by local marginal counts.
For a clause $f$, let $N_f((a_i)_{i\in I_f})$ count its ordered triples
of satisfying assignments having exactly $a_i$ ones at variable $i$.
In $Q_i^3$, choosing the zero-labelled monomial $a_i$ times has
coefficient $\binom3{a_i}$. Balance at each incidence forces the same
ones count $a_i$ in its clause triples. Therefore
\[
 A_2(\varphi)=
 \sum_{0\le a_i\le3}
   \prod_i\binom3{a_i}
   \prod_fN_f((a_i)_{i\in I_f}).
\]

For $\varphi_0$, the unit clauses force $a_1=a_2=3$. The required
tautological unary relation for the unused variable $x_0$ contributes
$\binom3{a_0}$, in addition to the same factor from $Q_0^3$. Hence
\[
 A_2(\varphi_0)=\sum_{a=0}^3\binom3a^2
 =1+9+9+1=20.
\]

For $\varphi_1$, put $a=a_0$, $b=a_2$, and again $a_1=3$.
Two clause relations are $\neg x_0\vee x_2$ and one is $x_0\vee x_2$.
Their ordered-triple counts are
\[
 N_-(a,b)=\binom3b\binom ba,\qquad
 N_+(a,b)=\binom3b\binom b{a+b-3},
\]
with out-of-range binomial coefficients zero. The first formula chooses
the $b$ positions of $x_2=1$ and the $a$ positions of $x_0=1$ within
them. For the second, $x_0$ must equal one wherever $x_2$ is zero.
Thus
\[
 A_2(\varphi_1)=
 \sum_{a,b=0}^3\binom3a\binom3b\,N_-(a,b)^2N_+(a,b).
\]
Exactly six summands are nonzero:
\[
 \begin{array}{c|r}
 (a,b)&\binom3a\binom3b\,N_-(a,b)^2N_+(a,b)\\ \hline
 (1,2)&972\\
 (2,2)&486\\
 (0,3)&1\\
 (1,3)&81\\
 (2,3)&81\\
 (3,3)&1
 \end{array}
\]
Their sum is $1622$. In particular,
\[
 \mathcal C(\varphi_0)=8,\qquad
 \mathcal C(\varphi_1)=809.
\]
Theorem~\ref{fg:thm:orbit} and \eqref{fg:eq:exact-height} give height
at least three for the first constructed group and height exactly two
for the second. The second invariant is therefore not determined by
$S$, or even by the Boolean function.

The example can use a common ambient space. The incidence counts are
three and nine, so the tagged input dimensions are seven and nineteen.
The common choice $Q=32$ is admissible for both. Their stabilized
equations then both have degree $33$ in $\PP^{32}$.
\end{example}

\subsection{Counting coefficient digits for arbitrary circuits}
\begin{lemma}[Fixed prime-power tests]\label{fg:lem:digits}
Let $f$ be a $\#\mathrm P$ function, and fix an integer $a\ge1$.
The function
\[
 \Delta_a(f)=
 \prod_{j=0}^{a-1}
 \left(1+(p-1)\binom{f}{p^j}^{p-1}\right)
\]
belongs to $\#\mathrm P$ and satisfies
\[
 \Delta_a(f)\equiv
 \begin{cases}
 1,&p^a\mid f,\\
 0,&p^a\nmid f
 \end{cases}\pmod p.
\]
Fixed finite conjunctions and complements of these tests also have
such zero-or-one counting representatives modulo $p$.
\end{lemma}

\begin{proof}
Encode accepting paths by witness strings of one polynomial length,
fixing unused positions to zero. For a fixed integer $q\ge0$,
guessing $q$ accepting strings in strictly increasing lexicographic
order gives exactly $\binom{f}{q}$ accepting tuples. Thus this
binomial is a $\#\mathrm P$ function.

Write the nonnegative integer $f$ in base $p$ as
$f=\sum_j f_jp^j$. In characteristic $p$,
\[
 (1+T)^f=\prod_j(1+T^{p^j})^{f_j}.
\]
The lower-index factors have combined degree less than $p^j$, and
higher-index factors have no positive term of degree at most $p^j$.
Comparison at degree $p^j$ therefore gives
$\binom{f}{p^j}\equiv f_j\pmod p$.

For any accepting-path count $c$, the nonnegative count
$1+(p-1)c^{p-1}$ is one modulo $p$ when $c$ is zero, and zero
otherwise. Products and sums are realized by tuples of accepting
computations and disjoint labelled branches. Their number is fixed
here. Applying this zero test to all the first $a$ digits proves the
formula. Products of zero-or-one residues give conjunctions, and
$t\mapsto1+(p-1)t$ complements such a residue. No count division is
used.
\end{proof}

\begin{lemma}[Coefficient witnesses for circuit input]
\label{fg:lem:circuit-coeff}
Fix $J\ge1$. Suppose an arithmetic circuit over $\Fp$ is promised to
define a homogeneous degree-$d$ polynomial in $d\ge3$ variables whose
projective hypersurface is smooth and geometrically integral.
There is a deterministic polynomial-time construction of a circuit for a
homogeneous $p$-primitive integer lift $\widetilde F$ with nonnegative coefficients
such that each function
\[
 B_{p^j}(\widetilde F)
 =[(x_1\cdots x_d)^{p^j-1}]
       \widetilde F^{p^j-1},\qquad 1\le j\le J,
\]
is represented by a polynomial-time counting verifier. In particular
these are $\#\mathrm P$ functions of the circuit input.
\end{lemma}

\begin{proof}
First, the numerical dimension is polynomially bounded on the promise.
If some variable is absent, its coordinate point lies on the
hypersurface and every first derivative vanishes there, since $d\ge3$.
This would be a singular point. Thus every ambient variable occurs
syntactically. Even with binary coordinate labels, $d$ is at most the
number of distinct variable inputs, hence at most the input length.
A total verifier may reject larger invalid declarations before any
dimension-dependent construction.

Replace subtraction by addition and multiplication by $p-1$. If
needed, homogenize through degree $d$: retain degrees zero through $d$
at every gate, and use convolution of homogeneous components at a
multiplication gate. This uses polynomially many gates, without
expanding higher-degree components. Evaluate each degree-zero
subcircuit in $\Fp$ and replace it by its representative in
$\{0,\ldots,p-1\}$. Interpret the remaining component circuit over the
nonnegative integers and take its degree-$d$ output. This defines
$\widetilde F$ and reduces to $F$ modulo $p$. It is $p$-primitive, meaning that some coefficient is not divisible by $p$,
because $F$ is nonzero.

Folding the degree-zero subcircuits is necessary for the witness
bound. Repeated squaring of an integer constant could otherwise give
a double-exponentially large coefficient even in a circuit of bounded
positive degree. We are choosing a lift of the given polynomial over
$\Fp$, so such integer evaluations need not be preserved. Every
nonzero scalar leaf of the chosen lift is bounded by $p-1$.

Consider now a homogeneous nonnegative circuit with $M$ gates and
output degree $D$, with degree-zero subcircuits already replaced by
scalar leaves. A distributive parse selects one input at an addition
gate and both inputs at a multiplication gate. A contributing parse
has exactly $D$ variable leaves. Every multiplication with two
positive-degree inputs splits positive degree, so there are at most
$D-1$ such branchings. Between them, a path has length at most $M$.
Scalar children are leaves. Thus the entire unfolded parse has
$O(MD)$ vertices, despite sharing in the original circuit.

A nondeterministic verifier follows the parse in a fixed depth-first
order. At a scalar leaf $c$, it makes exactly $c$ accepting labelled
choices, rejecting unused binary labels; zero rejects. At variable
leaves it accumulates the exponent vector, and at the end it checks
the target monomial. This counts every distributive contribution with
its coefficient multiplicity, in polynomial time.

Apply the construction to $\widetilde F^{p^j-1}$. For fixed $p,J$,
its circuit size and degree $D=d(p^j-1)$ are polynomial in the input
length. The same chosen lift works for all $j\le J$.
Lemma~\ref{fg:lem:height-log} applies to that lift; no specified
integer evaluation of the original circuit is part of the input.
\end{proof}

\subsection{Fixed thresholds and exact-height seeds}
\begin{theorem}[Complexity of every fixed height threshold]
\label{fg:thm:threshold}
Fix a prime $p$ and an integer $r\ge2$. The following promise problem
is $\ModP$-complete under deterministic polynomial-time many-one
reductions: given a homogeneous arithmetic circuit over $\Fp$ for a
degree-$d$ polynomial in $d\ge3$ variables, promised to define a smooth
geometrically integral hypersurface $X\subset\PP^{d-1}$, decide whether
the top Artin--Mazur formal group has $h(X)\ge r$.
Infinite height is included in this predicate. The dimension is
allowed to grow, and hardness persists after imposing any fixed lower
bound on it.

For every fixed finite $h\ge1$, the exact-height predicate $h(X)=h$
on the same circuit promise belongs to $\ModP$.
\end{theorem}

\begin{proof}
Use the lift of Lemma~\ref{fg:lem:circuit-coeff}, taking $J=r-1$ for the
threshold predicate and $J=h$ for the exact-height predicate, and put
$f_j=B_{p^j}(\widetilde F)$. By \eqref{fg:eq:all-height},
height at least $r$ is equivalent to $p^j\mid f_j$ for $1\le j<r$.
Lemma~\ref{fg:lem:digits} supplies the accepting-path count
\[
 \prod_{j=1}^{r-1}\Delta_j(f_j),
\]
whose residue is one exactly on those inputs. Likewise
\[
 \left(\prod_{j=1}^{h-1}\Delta_j(f_j)\right)
       \left(1+(p-1)\Delta_h(f_h)\right)
\]
represents exact height $h$ by \eqref{fg:eq:exact-height}, with an
empty product equal to one. These total counting machines need only
agree on the promise; they do not recognize smoothness.

For hardness, start with a CNF $\varphi$ and take $a=r-1$ copies of
$G_\varphi$ in disjoint variable blocks. Their product
\[
 G^{(a)}=\prod_{\nu=1}^aG_\varphi(x^{(\nu)})
\]
has degree $au$ in $au$ variables. Apply Theorem~\ref{thm:stabilizer}
with a $p$-power $Q\ge\max\{au+2,p^{r-1}\}$.
It gives a smooth geometrically integral anticanonical hypersurface.
Coefficients separate across the blocks, and
Lemma~\ref{fg:lem:replica} yields
\[
 b_{p^j}(F_{G^{(a)}})=
       \bigl(C_jA_j(\varphi)\bigr)^a\qquad(1\le j<r).
\]
If $S\equiv0\pmod p$, every $A_j$ is divisible by $p$, so these
numerators have valuation at least $a=r-1\ge j$. The output height
is at least $r$. If $S\not\equiv0\pmod p$, its first numerator
$(C_1S)^a$ is a unit, so the height is one. Thus
\[
 h(F_{G^{(a)}})\ge r\iff S\equiv0\pmod p.
\]
The product input itself need not be smooth; its stabilization is the
promised output.

To give the reduction in the nonzero-residue convention for $\ModP$,
let $q$ be any source accepting-path count and form
$q'=1+(p-1)q^{p-1}$. This is a $\#\mathrm P$ count and
$q'\equiv0\pmod p$ exactly when $q\not\equiv0\pmod p$.
The parsimonious Boolean encoding used for Theorem~\ref{thm:main}
gives a width-three CNF with exactly $q'$ satisfying assignments.
Applying the preceding construction gives a many-one reduction in
the required direction.

At fixed $p,r$, the number of blocks and the lower bound $p^{r-1}$
are constants. The least admissible $p$-power $Q$ is less than $p$
times its numerical threshold, so the circuit size is polynomial.
Choosing $Q$ also to exceed any fixed dimension bound preserves
this cost.
\end{proof}

\begin{theorem}[Exact height from one fixed seed]\label{fg:thm:exact-seed}
Fix a prime $p$ and a finite integer $h\ge1$. If there exists a smooth
geometrically integral anticanonical hypersurface over $\Fp$ of
positive dimension whose top formal group has height $h$, then
deciding exact height $h$ in the growing-dimensional circuit class
of Theorem~\ref{fg:thm:threshold} is $\ModP$-complete. If no such
hypersurface exists, its promise yes-locus is empty.
\end{theorem}

\begin{proof}
Fix one seed equation $H$ of degree $e$ in $e\ge3$ variables and a
homogeneous $p$-primitive integer lift $\widetilde H$. Its equation is
fixed data, because $p,h$, and the seed are fixed. Write
$\beta_j=B_{p^j}(\widetilde H)$. Equations
\eqref{fg:eq:all-height}--\eqref{fg:eq:exact-height} give
\[
 p^j\mid\beta_j\quad(1\le j<h),\qquad
 v_p(\beta_h)=h-1.
\]

Given $\varphi$, form $G=\widetilde H(v)G_\varphi(w)$ in disjoint
variable blocks. It has degree $e+u$ in $e+u$ variables. Stabilize
its reduction modulo $p$ with a $p$-power
$Q\ge\max\{e+u+2,p^h\}$, using the displayed integer form for its
lift. Coefficient separation and
\eqref{fg:eq:stabilized-numerator} give
\[
 b_{p^j}(F_G)=\beta_jC_jA_j(\varphi)\qquad(1\le j\le h).
\]
All prior divisibilities hold for every input. At the $h$th stage,
$\beta_h$ has valuation $h-1$, $C_h$ is a unit, and
$A_h(\varphi)\equiv S^h\pmod p$. Consequently
\[
 h(F_G)=h\iff S\not\equiv0\pmod p.
\]
When $S\equiv0$, the output instead has height at least $h+1$,
possibly infinite. The map is polynomial time with the fixed seed
hard-coded, and every output satisfies the geometric promise.
This proves hardness; membership is
Theorem~\ref{fg:thm:threshold}.
\end{proof}

Theorem~\ref{fg:thm:all-heights} below supplies a seed over $\Fp$
for every finite height. Corollary~\ref{fg:cor:exact-height-complete}
therefore makes exact-height completeness unconditional.
The counting algorithms above treat $p$ and the finite threshold as
fixed parameters. They assert neither polynomial dependence on a
binary-encoded variable threshold nor a finite stopping rule for the
infinite list in \eqref{fg:eq:infinite-height}.
